% Unit 7 self-study -- "I do / You do" ladder for CED 7.1-7.12.
% Twelve ladders, forty-eight questions.
%
% Disjoint from u07-frq, which uses the PCl5 dissociation where the 5%
% approximation FAILS, a Q-vs-K item at K=0.0211, the NH4HS heterogeneous
% equilibrium, the isotope-labelling dynamic-equilibrium question and CaF2
% solubility. This document uses H2/I2, an approximation that SUCCEEDS,
% CaCO3 decomposition, and AgCl/PbI2/Ag2CrO4 solubility.
\documentclass{apchem-worksheet}

\apcunit{7}
\apcunittitle{Equilibrium}
\apcdoctitle{Self-Study \textbullet\ CED 7.1--7.12, I~do / You~do}
\apcsubtitle{Twelve ladders \textbullet\ four YOUR TURN questions each \textbullet\ work alone, check after all four}

\begin{document}
\apctitleblock

\begin{apcnote}[How to use these notes --- read this first]
Each skill is a \textbf{ladder}: a worked example, then \textbf{four}
YOUR TURN questions --- same skill, new numbers, no help. Work all four
before checking anything.

\textbf{The distinction this whole unit rests on.} $K$ is a
\textbf{constant} at a given temperature; $Q$ is the \textbf{same
expression evaluated right now}, at whatever concentrations happen to
exist. Comparing the two tells you which way the reaction will move. Every
Le Ch\^{a}telier question is really a $Q$-versus-$K$ question underneath.

\textbf{The habit that decides this unit.} Only \emph{temperature} changes
the value of $K$. Adding reactant, removing product, changing the volume
or adding a catalyst all leave $K$ exactly where it was --- they move the
system along a fixed $K$, they do not move $K$.

\textbf{One scope note.} Converting between $K_c$ and $K_p$ is
\textbf{excluded} by the CED. The conceptual difference between them is
examinable; the $K_p = K_c(RT)^{\Delta n}$ conversion is not, and does not
appear here.
\end{apcnote}

% =====================================================================
\section*{\sffamily Ladder 1 \textbullet\ What equilibrium is}
\ced{7.1} \ced{7.2}

Equilibrium is \emph{dynamic}: both reactions continue, at equal rates.
Concentrations stop changing; the molecules do not stop reacting.

\begin{workedexample}[I do: what stops and what does not]
Describe what is happening at equilibrium, macroscopically and at the
particle level.

\textbf{Macroscopically: nothing changes.} Concentrations, colour and
pressure are all constant. On a concentration-versus-time graph, every
curve has flattened.

\textbf{At the particle level: everything continues.} The forward reaction
is still converting reactants to products, and the reverse is still
converting them back --- at exactly the \emph{same rate}. Every molecule
consumed one way is replaced by one made the other way.

\textbf{Why the graph cannot prove this.} If the reaction had simply
stopped, the curves would also flatten. The evidence for dynamic behaviour
comes from experiments such as isotopic labelling, where a labelled
molecule placed on one side later appears on both.

\textbf{Also note: reactants are still present.} A reaction at equilibrium
has not gone to completion, and never will --- which is why the arrow is
drawn both ways.
\end{workedexample}

\begin{yourturn}[YOUR TURN 1 --- four questions]
\begin{enumerate}[leftmargin=*,itemsep=0.5em]
  \item At equilibrium, what is true of the forward and reverse rates?
        \workspace[1.4cm]
  \item Are reactants still present at equilibrium? \workspace[1.4cm]
  \item What does a concentration-time graph look like once equilibrium is
        reached? \workspace[1.6cm]
  \item Why does that graph not prove the process is dynamic?
        \workspace[1.8cm]
\end{enumerate}
\selfcheck{(a) equal \quad (b) yes \quad (c) horizontal lines at non-zero
values \quad (d) a stopped reaction would look identical}
\keyonly{\begin{apcanswer}(a) They are \textbf{equal}. That equality is
what makes the concentrations constant.
(b) \textbf{Yes.} Equilibrium is not completion; both reactants and
products are present in fixed amounts.
(c) All curves become \textbf{horizontal} at non-zero concentrations ---
constant, but not zero.
(d) Because a reaction that had genuinely \emph{stopped} would produce
exactly the same flat lines. The graph shows only that nothing is changing
macroscopically. Distinguishing the two requires evidence at the molecular
level, such as isotopic labelling.\end{apcanswer}}
\end{yourturn}

% =====================================================================
\section*{\sffamily Ladder 2 \textbullet\ Writing equilibrium expressions}
\ced{7.3}

Products over reactants, each raised to its coefficient. Pure solids and
pure liquids are omitted.

\begin{workedexample}[I do: two expressions, one omission]
Write $K$ for
\ce{N2(g) + 3H2(g) <=> 2NH3(g)} and for
\ce{CaCO3(s) <=> CaO(s) + CO2(g)}.

\textbf{First:} all three are gases, so all appear, each raised to its
coefficient.
\[ K = \frac{[\ce{NH3}]^2}{[\ce{N2}][\ce{H2}]^3} \]

\textbf{Second:} both \ce{CaCO3} and \ce{CaO} are \textbf{pure solids} and
are omitted.
\[ K = [\ce{CO2}] \]

\textbf{Why solids drop out.} A pure solid has a fixed density and
composition, so its ``concentration'' does not change no matter how much
is present. Adding more solid cannot shift the equilibrium, so it is
absorbed into the constant rather than written.

\textbf{The consequence worth noticing:} in the second reaction, the
pressure of \ce{CO2} above the solid at a given temperature is fixed by
$K$ alone --- however much or little \ce{CaCO3} is in the flask.
\end{workedexample}

\begin{yourturn}[YOUR TURN 2 --- four questions]
\begin{enumerate}[leftmargin=*,itemsep=0.5em]
  \item Write $K$ for \ce{2SO2(g) + O2(g) <=> 2SO3(g)}:
        \workspace[1.6cm]
  \item Write $K$ for \ce{H2O(l) <=> H2O(g)}: \workspace[1.4cm]
  \item Write $K$ for \ce{Fe3O4(s) + 4H2(g) <=> 3Fe(s) + 4H2O(g)}:
        \workspace[1.6cm]
  \item Why does adding more solid not shift such an equilibrium?
        \workspace[1.8cm]
\end{enumerate}
\selfcheck{(a) $[\ce{SO3}]^2/([\ce{SO2}]^2[\ce{O2}])$ \quad
(b) $K = [\ce{H2O}(g)]$ \quad (c) $[\ce{H2O}]^4/[\ce{H2}]^4$ \quad
(d) its concentration is constant}
\keyonly{\begin{apcanswer}(a) $K = \dfrac{[\ce{SO3}]^2}
{[\ce{SO2}]^2[\ce{O2}]}$.
(b) $K = [\ce{H2O}(g)]$ --- the liquid is omitted. This constant is just
the vapour pressure of water at that temperature.
(c) $K = \dfrac{[\ce{H2O}]^4}{[\ce{H2}]^4}$ --- both solids omitted.
(d) Because a pure solid's concentration is fixed by its density and does
not appear in $K$. Adding more increases the \emph{amount} but not the
concentration, so there is nothing for the system to respond
to.\end{apcanswer}}
\end{yourturn}

% =====================================================================
\section*{\sffamily Ladder 3 \textbullet\ Calculating K from equilibrium concentrations}
\ced{7.4}

If you are handed the equilibrium concentrations, substitute them
straight in. The only difficulty is remembering the exponents.

\begin{workedexample}[I do: substitution]
At equilibrium, \ce{H2 + I2 <=> 2HI} has
$[\ce{HI}] = \SI{0.50}{\Molar}$ and
$[\ce{H2}] = [\ce{I2}] = \SI{0.10}{\Molar}$. Find $K$.

\[ K = \frac{[\ce{HI}]^2}{[\ce{H2}][\ce{I2}]}
     = \frac{(0.50)^2}{(0.10)(0.10)}
     = \frac{0.25}{0.010} = \mathbf{25} \]

\textbf{The exponent is the whole difficulty.} Forgetting to square
\ce{HI} gives \num{50} --- exactly double, and a common wrong answer.

\textbf{Sense check:} $K = 25$ is comfortably greater than 1, so products
should dominate at equilibrium. Indeed \SI{0.50}{\Molar} of \ce{HI}
against \SI{0.10}{\Molar} of each reactant. \checkmark

\textbf{$K$ has no units on this exam.} Report it as a bare number.
\end{workedexample}

\begin{yourturn}[YOUR TURN 3 --- four questions]
\begin{enumerate}[leftmargin=*,itemsep=0.5em]
  \item $[\ce{NH3}] = 0.40$, $[\ce{N2}] = 0.20$, $[\ce{H2}] = 0.20$ for
        \ce{N2 + 3H2 <=> 2NH3}. Find $K$. \workspace[1.8cm]
  \item For \ce{2NO2 <=> N2O4} with $[\ce{N2O4}] = 0.60$ and
        $[\ce{NO2}] = 0.20$, find $K$. \workspace[1.6cm]
  \item If $K$ for a reaction is 4, what is $K$ for its reverse?
        \workspace[1.4cm]
  \item Why is the coefficient used as an exponent rather than a
        multiplier? \workspace[1.8cm]
\end{enumerate}
\selfcheck{(a) 100 \quad (b) 15 \quad (c) $1/4 = 0.25$ \quad (d) $K$
derives from a product of concentrations}
\keyonly{\begin{apcanswer}(a) $K = (0.40)^2/[(0.20)(0.20)^3] =
0.16/0.0016 = \mathbf{100}$.
(b) $K = 0.60/(0.20)^2 = 0.60/0.040 = \mathbf{15}$.
(c) $\mathbf{0.25}$. Reversing a reaction inverts $K$, since numerator and
denominator swap.
(d) Because the equilibrium expression comes from a \emph{product} of
concentration terms, one for each particle taking part. Two \ce{NO2}
molecules contribute $[\ce{NO2}] \times [\ce{NO2}]$, which is
$[\ce{NO2}]^2$ --- multiplication, not addition.\end{apcanswer}}
\end{yourturn}

% =====================================================================
\section*{\sffamily Ladder 4 \textbullet\ What the magnitude of K tells you}
\ced{7.5}

$K$ measures \emph{extent} --- how far the reaction proceeds --- and
nothing else. It says nothing about rate.

\begin{workedexample}[I do: reading three values]
Interpret $K = 1\times10^{10}$, $K = 1$ and
$K = 1\times10^{-8}$.

\textbf{$K = 10^{10}$:} equilibrium lies far to the \textbf{right}. At
equilibrium the mixture is essentially all products, with only traces of
reactant. Such a reaction is often said to ``go to completion'', though
strictly a little reactant always remains.

\textbf{$K = 1$:} neither side dominates --- comparable amounts of both are
present.

\textbf{$K = 10^{-8}$:} far to the \textbf{left}. Almost nothing reacts;
the mixture is essentially all reactant.

\textbf{What none of these tells you: the rate.} A reaction with
$K = 10^{10}$ may be immeasurably slow if its activation energy is large
--- the conversion of diamond to graphite is exactly that. Extent and rate
are independent properties.
\end{workedexample}

\begin{yourturn}[YOUR TURN 4 --- four questions]
\begin{enumerate}[leftmargin=*,itemsep=0.5em]
  \item $K = 5\times10^{6}$. Which side is favoured?
        \workspace[1.4cm]
  \item $K = 3\times10^{-5}$. Which side? \workspace[1.4cm]
  \item Does a large $K$ mean a fast reaction? \workspace[1.6cm]
  \item Two reactions have $K = 100$ and $K = 0.01$. How are they
        related? \workspace[1.6cm]
\end{enumerate}
\selfcheck{(a) products \quad (b) reactants \quad (c) no \quad (d) they
are the reverse of each other}
\keyonly{\begin{apcanswer}(a) \textbf{Products} --- $K \gg 1$.
(b) \textbf{Reactants} --- $K \ll 1$.
(c) \textbf{No.} $K$ describes only how far the reaction goes, not how
quickly. Rate depends on activation energy, which is entirely separate.
(d) They are \textbf{reverses} of one another: $0.01 = 1/100$. Inverting a
reaction inverts its equilibrium constant.\end{apcanswer}}
\end{yourturn}

% =====================================================================
\section*{\sffamily Ladder 5 \textbullet\ Q versus K}
\ced{7.3} \ced{7.8}

$Q$ has exactly the same form as $K$ but is evaluated at the
\emph{current} concentrations. Comparing them predicts the direction of
change.

\begin{workedexample}[I do: predicting the shift]
For \ce{H2 + I2 <=> 2HI} with $K = 25$, a mixture holds
$[\ce{HI}] = 0.20$, $[\ce{H2}] = 0.30$,
$[\ce{I2}] = 0.30$~M. Which way does it move?

\[ Q = \frac{(0.20)^2}{(0.30)(0.30)} = \frac{0.040}{0.090}
   = \mathbf{0.44} \]

\textbf{Compare:} $Q = 0.44$ is much less than $K = 25$. There is too
little product relative to equilibrium, so the reaction proceeds
\textbf{forward}, to the right, making more \ce{HI} until $Q$ climbs to
25.

\textbf{The three cases, once and for all:}
$Q < K$ shift right (make products);
$Q = K$ already at equilibrium, no net change;
$Q > K$ shift left (make reactants).

\textbf{Why this is the master tool:} every Le Ch\^{a}telier prediction is
this comparison in disguise. Adding a reactant lowers $Q$; the system
responds by moving right. Nothing else is going on.
\end{workedexample}

\begin{yourturn}[YOUR TURN 5 --- four questions]
\begin{enumerate}[leftmargin=*,itemsep=0.5em]
  \item $Q = 0.5$, $K = 2.0$. Which way? \workspace[1.4cm]
  \item $Q = 150$, $K = 25$. Which way? \workspace[1.4cm]
  \item $Q = K$. What happens? \workspace[1.4cm]
  \item Adding more reactant to a mixture at equilibrium: what happens to
        $Q$, and then to the system? \workspace[1.8cm]
\end{enumerate}
\selfcheck{(a) right \quad (b) left \quad (c) nothing --- already at
equilibrium \quad (d) $Q$ falls; the system shifts right}
\keyonly{\begin{apcanswer}(a) \textbf{Right} --- $Q < K$, too little
product.
(b) \textbf{Left} --- $Q > K$, too much product.
(c) \textbf{Nothing.} The system is already at equilibrium and there is no
net change in either direction.
(d) Reactant appears in the \emph{denominator} of $Q$, so increasing it
makes $Q$ \textbf{smaller}. Now $Q < K$, and the system shifts
\textbf{right} to consume some of the added reactant and restore
$Q = K$.\end{apcanswer}}
\end{yourturn}

% =====================================================================
\section*{\sffamily Ladder 6 \textbullet\ ICE tables}
\ced{7.7}

Initial, Change, Equilibrium. The change row carries the
\textbf{coefficients} as multipliers of $x$ --- that is where most errors
enter.

\begin{workedexample}[I do: setting one up correctly]
\SI{0.500}{\Molar} \ce{PCl5} is placed in a flask and dissociates by
\ce{PCl5 <=> PCl3 + Cl2}. Build the ICE table.

\begin{center}
\begin{tabular}{lccc}
\hline
 & \ce{PCl5} & \ce{PCl3} & \ce{Cl2} \\
\hline
Initial     & 0.500     & 0        & 0 \\
Change      & $-x$      & $+x$     & $+x$ \\
Equilibrium & $0.500-x$ & $x$      & $x$ \\
\hline
\end{tabular}
\end{center}

\textbf{Here every coefficient is 1}, so every change is $\pm x$.

\textbf{But if the equation were \ce{H2 + I2 <=> 2HI}}, the change row
would read $-x$, $-x$, $+\mathbf{2}x$ --- the coefficient 2 multiplies.
Writing $+x$ for \ce{HI} is the single commonest ICE error and corrupts
every later line.

\textbf{Reading the table:} the equilibrium row is what goes into $K$.
Nothing else does.
\end{workedexample}

\begin{yourturn}[YOUR TURN 6 --- four questions]
\begin{enumerate}[leftmargin=*,itemsep=0.5em]
  \item For \ce{N2 + 3H2 <=> 2NH3} starting from \SI{1.0}{\Molar} of each
        reactant, write the change row. \workspace[1.6cm]
  \item For \ce{2NO2 <=> N2O4} starting from \SI{0.80}{\Molar} \ce{NO2},
        write the equilibrium row. \workspace[1.6cm]
  \item Which row is substituted into the $K$ expression?
        \workspace[1.4cm]
  \item Why must coefficients appear in the change row?
        \workspace[1.8cm]
\end{enumerate}
\selfcheck{(a) $-x$, $-3x$, $+2x$ \quad (b) $0.80 - 2x$ and $x$ \quad
(c) the equilibrium row \quad (d) species react in fixed proportions}
\keyonly{\begin{apcanswer}(a) $-x$, $-3x$, $+2x$.
(b) \ce{NO2}: $0.80 - 2x$; \ce{N2O4}: $x$.
(c) The \textbf{equilibrium} row.
(d) Because the balanced equation fixes the proportions in which species
are consumed and produced. If $x$ moles of \ce{N2} react, exactly $3x$ of
\ce{H2} must go with it and $2x$ of \ce{NH3} must appear. Using $x$
everywhere would describe a reaction with different
stoichiometry.\end{apcanswer}}
\end{yourturn}

% =====================================================================
\section*{\sffamily Ladder 7 \textbullet\ The 5\% approximation}
\ced{7.7}

When $K$ is very small, $x$ is tiny compared with the initial
concentration and the subtraction can be dropped. \textbf{Always test the
assumption afterwards.}

\begin{workedexample}[I do: assume, solve, then check]
For \ce{A <=> B + C} with $K = 1.0\times10^{-5}$ and
$[\ce{A}]_0 = \SI{0.100}{\Molar}$, find $x$.

\textbf{Set up:}
$\dfrac{x^2}{0.100 - x} = 1.0\times10^{-5}$.

\textbf{Assume $x \ll 0.100$}, so the denominator is about \num{0.100}:
\[ x = \sqrt{(1.0\times10^{-5})(0.100)} = \mathbf{1.0\times10^{-3}} \]

\textbf{Now test it.} As a fraction of the initial concentration,
$1.0\times10^{-3}/0.100 = \SI{1.0}{\percent}$ --- comfortably below the
\SI{5}{\percent} threshold, so the assumption is \textbf{valid} and no
quadratic is needed.

\textbf{What ``valid'' means:} solving the quadratic properly gives
$x = 9.95\times10^{-4}$, differing by \SI{0.5}{\percent} --- far less than
the precision of the data.

\textbf{When it fails,} as it does whenever $K$ is not small, you must use
the quadratic. Testing takes ten seconds; skipping the test is what costs
marks.
\end{workedexample}

\begin{yourturn}[YOUR TURN 7 --- four questions]
\begin{enumerate}[leftmargin=*,itemsep=0.5em]
  \item What is the usual threshold for accepting the approximation?
        \workspace[1.4cm]
  \item $x = 0.004$ from an initial \SI{0.200}{\Molar}. Is it valid?
        \workspace[1.6cm]
  \item $x = 0.03$ from an initial \SI{0.10}{\Molar}. Is it valid?
        \workspace[1.6cm]
  \item For which kind of $K$ is the approximation generally safe?
        \workspace[1.6cm]
\end{enumerate}
\selfcheck{(a) \SI{5}{\percent} \quad (b) yes --- \SI{2}{\percent} \quad
(c) no --- \SI{30}{\percent} \quad (d) very small $K$}
\keyonly{\begin{apcanswer}(a) $x$ must be under \textbf{\SI{5}{\percent}}
of the initial concentration.
(b) $0.004/0.200 = \SI{2.0}{\percent}$ --- \textbf{valid}.
(c) $0.03/0.10 = \SI{30}{\percent}$ --- \textbf{not valid}; the quadratic
is required.
(d) When $K$ is \textbf{very small}. A small $K$ means the reaction barely
proceeds, so $x$ stays tiny relative to what was there initially. When $K$
is large the reaction goes far, $x$ is a big fraction of the initial
amount, and the approximation collapses.\end{apcanswer}}
\end{yourturn}

% =====================================================================
\section*{\sffamily Ladder 8 \textbullet\ Le Ch\^{a}telier: concentration}
\ced{7.9} \ced{7.10}

Add something and the system consumes it; remove something and the system
replaces it. $K$ does not change.

\begin{workedexample}[I do: four changes to one equilibrium]
For \ce{N2 + 3H2 <=> 2NH3}, state the effect of: adding \ce{N2}; removing
\ce{NH3}; adding \ce{NH3}; adding argon at constant volume.

\textbf{Adding \ce{N2}:} shifts \textbf{right}. The system consumes some
of the added reactant. ($Q$ falls below $K$.)

\textbf{Removing \ce{NH3}:} shifts \textbf{right}. The system replaces
what was taken. ($Q$ falls.)

\textbf{Adding \ce{NH3}:} shifts \textbf{left}. ($Q$ rises above $K$.)

\textbf{Adding argon at constant volume:} \textbf{no shift}. Argon takes
no part in the reaction and does not change the concentration of anything
that does. The total pressure rises, but every partial pressure that
matters is unchanged, so $Q$ is unchanged.

\textbf{Throughout, $K$ is constant} --- only the position along it moved.
\end{workedexample}

\begin{yourturn}[YOUR TURN 8 --- four questions]
\begin{enumerate}[leftmargin=*,itemsep=0.5em]
  \item For \ce{2SO2 + O2 <=> 2SO3}, effect of adding \ce{O2}?
        \workspace[1.4cm]
  \item Effect of removing \ce{SO3}? \workspace[1.4cm]
  \item Effect of adding helium at constant volume?
        \workspace[1.4cm]
  \item Does adding a reactant change $K$? \workspace[1.4cm]
\end{enumerate}
\selfcheck{(a) right \quad (b) right \quad (c) none \quad (d) no}
\keyonly{\begin{apcanswer}(a) Shifts \textbf{right}, consuming some of the
added oxygen.
(b) Shifts \textbf{right}, replacing the removed product.
(c) \textbf{No effect.} Helium is inert and, at constant volume, does not
change the concentration or partial pressure of any participating species.
(d) \textbf{No.} $K$ depends only on temperature. Adding a reactant moves
the system to a new position along the \emph{same} $K$.\end{apcanswer}}
\end{yourturn}

% =====================================================================
\section*{\sffamily Ladder 9 \textbullet\ Le Ch\^{a}telier: volume and pressure}
\ced{7.9} \ced{7.10}

Compressing a gaseous equilibrium shifts it toward the side with
\emph{fewer} moles of gas. Count the gas coefficients on each side ---
solids and liquids do not count.

\begin{workedexample}[I do: counting moles of gas]
Predict the effect of halving the volume for:
(i) \ce{N2 + 3H2 <=> 2NH3};
(ii) \ce{H2 + I2 <=> 2HI}.

\textbf{(i)} Left has $1 + 3 = 4$ moles of gas; right has 2. Compressing
raises all pressures, and the system relieves that by moving to the side
with fewer particles: \textbf{shifts right}.

\textbf{(ii)} Left has $1 + 1 = 2$; right has 2. \textbf{Equal}, so
compression has \textbf{no effect} on the position --- there is no side
with fewer moles to move toward.

\textbf{The rule in one line:} equal moles of gas on both sides means
volume changes do nothing.

\textbf{And note what a volume change never does:} it does not alter $K$.
Only temperature does that.
\end{workedexample}

\begin{yourturn}[YOUR TURN 9 --- four questions]
\begin{enumerate}[leftmargin=*,itemsep=0.5em]
  \item Effect of compressing \ce{2NO2 <=> N2O4}?
        \workspace[1.4cm]
  \item Effect of expanding that same system? \workspace[1.4cm]
  \item Effect of compressing \ce{CO(g) + H2O(g) <=> CO2(g) + H2(g)}?
        \workspace[1.6cm]
  \item In \ce{CaCO3(s) <=> CaO(s) + CO2(g)}, how many moles of gas are on
        each side? \workspace[1.6cm]
\end{enumerate}
\selfcheck{(a) right, toward 1 mole \quad (b) left \quad (c) none ---
2 vs 2 \quad (d) 0 left, 1 right}
\keyonly{\begin{apcanswer}(a) Shifts \textbf{right}: 2 moles of gas become
1, so the product side relieves the pressure.
(b) Shifts \textbf{left}, toward the side with \emph{more} gas.
(c) \textbf{No effect} --- two moles of gas on each side.
(d) \textbf{Zero} on the left (\ce{CaCO3} is a solid) and \textbf{one} on
the right. Compressing therefore shifts it left, toward the solids.
Solids are excluded from the count for exactly the same reason they are
excluded from $K$.\end{apcanswer}}
\end{yourturn}

% =====================================================================
\section*{\sffamily Ladder 10 \textbullet\ Le Ch\^{a}telier: temperature and catalysts}
\ced{7.9}

Temperature is the \emph{only} change that alters $K$. A catalyst alters
neither $K$ nor the position.

\begin{workedexample}[I do: heat as a participant]
For the exothermic \ce{2SO2 + O2 <=> 2SO3}, $\Delta H < 0$, state the
effect of raising the temperature, and of adding a catalyst.

\textbf{Treat heat as a product} for an exothermic reaction:
\ce{2SO2 + O2 <=> 2SO3 + heat}.

\textbf{Raising the temperature} adds heat, so the system shifts to
consume it --- to the \textbf{left}. Less \ce{SO3} at equilibrium.

\textbf{And $K$ itself decreases.} This is the one case where the constant
genuinely changes value, because $K$ is defined at a temperature.

\textbf{Adding a catalyst: nothing changes.} It lowers the activation
energy of the forward and reverse reactions equally, so both rates rise in
the same proportion. Equilibrium arrives sooner and lies in exactly the
same place.

\textbf{The industrial tension this creates:} a lower temperature would
give more \ce{SO3}, but too slowly to be useful. The Contact process
compromises --- moderate temperature plus a catalyst to make that
temperature fast enough.
\end{workedexample}

\begin{yourturn}[YOUR TURN 10 --- four questions]
\begin{enumerate}[leftmargin=*,itemsep=0.5em]
  \item For an endothermic reaction, effect of raising the temperature?
        \workspace[1.6cm]
  \item Which change alters the value of $K$? \workspace[1.4cm]
  \item Effect of a catalyst on the equilibrium yield?
        \workspace[1.4cm]
  \item Why does a catalyst not shift the position?
        \workspace[1.8cm]
\end{enumerate}
\selfcheck{(a) shifts right; $K$ increases \quad (b) temperature only
\quad (c) none \quad (d) it speeds both directions equally}
\keyonly{\begin{apcanswer}(a) Shifts \textbf{right}. For an endothermic
reaction heat behaves as a \emph{reactant}, so supplying it drives the
reaction forward --- and $K$ \textbf{increases}.
(b) \textbf{Temperature}, and only temperature.
(c) \textbf{None.} The yield at equilibrium is unchanged; only the time
taken to reach it is shorter.
(d) Because it lowers the activation energy for the forward and reverse
reactions by the \emph{same} amount --- both directions cross the same
transition state. Both rates rise in the same proportion, so the point at
which they become equal is exactly where it was.\end{apcanswer}}
\end{yourturn}

% =====================================================================
\section*{\sffamily Ladder 11 \textbullet\ Solubility equilibria and \emph{K}\textsubscript{sp}}
\ced{7.11}

$K_{sp}$ is just an equilibrium constant for a dissolving solid. The solid
is omitted, and the stoichiometric coefficients become exponents --- and
multipliers of $s$.

\begin{workedexample}[I do: two different stoichiometries]
Find the molar solubility of \ce{AgCl}
($K_{sp} = 1.8\times10^{-10}$) and of \ce{PbI2}
($K_{sp} = 7.1\times10^{-9}$).

\textbf{\ce{AgCl}:} \ce{AgCl(s) <=> Ag+ + Cl^-}, a $1:1$ dissolution, so
$[\ce{Ag+}] = [\ce{Cl-}] = s$.
\[ K_{sp} = s^2 \;\Rightarrow\; s = \sqrt{1.8\times10^{-10}}
   = \mathbf{1.3\times10^{-5}~M} \]

\textbf{\ce{PbI2}:} \ce{PbI2(s) <=> Pb^2+ + 2I^-}, so
$[\ce{Pb^2+}] = s$ but $[\ce{I-}] = \mathbf{2}s$.
\[ K_{sp} = (s)(2s)^2 = 4s^3 = 7.1\times10^{-9} \]
\[ s^3 = 1.775\times10^{-9} \;\Rightarrow\;
   s = \mathbf{1.2\times10^{-3}~M} \]

\textbf{The factor of 2 is the entire difficulty.} Writing
$K_{sp} = s^3$ instead of $4s^3$ gives \num{1.9e-3} --- wrong, and a very
common answer.

\textbf{A warning about comparing:} \ce{PbI2} has the smaller $K_{sp}$ yet
is nearly a hundred times \emph{more} soluble than \ce{AgCl}. $K_{sp}$
values may only be compared directly between salts of the same
stoichiometric type.
\end{workedexample}

\begin{yourturn}[YOUR TURN 11 --- four questions]
\begin{enumerate}[leftmargin=*,itemsep=0.5em]
  \item Write the $K_{sp}$ expression for \ce{Ag2CrO4}:
        \workspace[1.6cm]
  \item Its $K_{sp}$ is $1.1\times10^{-12}$. Find its molar solubility.
        \workspace[2.0cm]
  \item In that saturated solution, what is $[\ce{Ag+}]$?
        \workspace[1.4cm]
  \item Why can $K_{sp}$ values not always be compared directly?
        \workspace[1.8cm]
\end{enumerate}
\selfcheck{(a) $[\ce{Ag+}]^2[\ce{CrO4^2-}]$ \quad (b)
\SI{6.5e-5}{\Molar} \quad (c) \SI{1.3e-4}{\Molar} \quad (d) different
stoichiometries}
\keyonly{\begin{apcanswer}(a) $K_{sp} = [\ce{Ag+}]^2[\ce{CrO4^2-}]$.
(b) $[\ce{Ag+}] = 2s$, $[\ce{CrO4^2-}] = s$, so
$K_{sp} = (2s)^2(s) = 4s^3 = 1.1\times10^{-12}$;
$s^3 = 2.75\times10^{-13}$, giving
$s = \mathbf{6.5\times10^{-5}}$~M.
(c) $2s = \mathbf{1.3\times10^{-4}}$~M.
(d) Because the relationship between $K_{sp}$ and solubility depends on
the dissolution stoichiometry. For an \ce{MX} salt $K_{sp} = s^2$, for
\ce{MX2} it is $4s^3$. A smaller $K_{sp}$ of a different type can still
correspond to a larger solubility, as \ce{PbI2} against \ce{AgCl}
shows.\end{apcanswer}}
\end{yourturn}

% =====================================================================
\section*{\sffamily Ladder 12 \textbullet\ The common-ion effect}
\ced{7.12}

Adding an ion that already appears in the equilibrium pushes it backwards.
It is Le Ch\^{a}telier applied to dissolving.

\begin{workedexample}[I do: how much less soluble?]
\ce{AgCl} has $K_{sp} = 1.8\times10^{-10}$. Compare its solubility in
pure water with that in \SI{0.10}{\Molar} \ce{NaCl}.

\textbf{In pure water:}
$s = \sqrt{1.8\times10^{-10}} = \SI{1.3e-5}{\Molar}$.

\textbf{In \SI{0.10}{\Molar} \ce{NaCl}:} chloride is already present at
\SI{0.10}{\Molar}, overwhelmingly more than the dissolution could supply,
so $[\ce{Cl-}] \approx 0.10$.
\[ K_{sp} = [\ce{Ag+}](0.10) \;\Rightarrow\;
   [\ce{Ag+}] = s = \frac{1.8\times10^{-10}}{0.10}
   = \mathbf{1.8\times10^{-9}~M} \]

\textbf{That is about 7000 times less soluble.} The common ion suppresses
dissolution dramatically.

\textbf{And $K_{sp}$ is unchanged throughout.} What changed is how the
fixed product of concentrations is distributed between the two ions ---
much more chloride, so far less silver.
\end{workedexample}

\begin{yourturn}[YOUR TURN 12 --- four questions]
\begin{enumerate}[leftmargin=*,itemsep=0.5em]
  \item Is \ce{AgCl} more or less soluble in \SI{0.10}{\Molar}
        \ce{AgNO3} than in water? \workspace[1.6cm]
  \item Does $K_{sp}$ change when a common ion is added?
        \workspace[1.4cm]
  \item Which direction does the dissolution equilibrium shift?
        \workspace[1.4cm]
  \item Would adding \ce{NaNO3} change \ce{AgCl}'s solubility?
        \workspace[1.8cm]
\end{enumerate}
\selfcheck{(a) less soluble \quad (b) no \quad (c) left \quad (d) no ---
neither ion is common}
\keyonly{\begin{apcanswer}(a) \textbf{Less soluble.} \ce{Ag+} is a common
ion, so the equilibrium is pushed back toward the solid.
(b) \textbf{No.} $K_{sp}$ depends only on temperature.
(c) To the \textbf{left}, toward undissolved solid.
(d) \textbf{No} (to a first approximation). Neither \ce{Na+} nor
\ce{NO3-} appears in the \ce{AgCl} dissolution equilibrium, so neither is
a common ion and $Q$ is unaffected. Only ions that appear in the
equilibrium expression can shift it.\end{apcanswer}}
\end{yourturn}

% =====================================================================
\section*{\sffamily Where you stand}

Tick a ladder only if all four YOUR TURN questions were right first time.

\renewcommand{\arraystretch}{1.30}
\noindent\begin{tabular}{@{}c p{6.0cm} c >{\raggedright\arraybackslash}p{4.9cm}@{}}
\toprule
\textbf{First try?} & \textbf{Skill} & \textbf{Ladder} &
  \textbf{If not, re-read\ldots}\\
\midrule
$\square$ & what equilibrium is & 1 & dynamic, not stopped \\
$\square$ & writing $K$ & 2 & omit solids and liquids \\
$\square$ & calculating $K$ & 3 & coefficients are exponents \\
$\square$ & magnitude of $K$ & 4 & extent, never rate \\
$\square$ & $Q$ versus $K$ & 5 & $Q<K$ right, $Q>K$ left \\
$\square$ & ICE tables & 6 & coefficients multiply $x$ \\
$\square$ & the 5\% test & 7 & always check afterwards \\
$\square$ & Le Ch.: concentration & 8 & consume what is added \\
$\square$ & Le Ch.: volume & 9 & count moles of gas \\
$\square$ & Le Ch.: temperature & 10 & only this changes $K$ \\
$\square$ & $K_{sp}$ and solubility & 11 & watch the $2s$ \\
$\square$ & common-ion effect & 12 & shifts left; $K_{sp}$ fixed \\
\bottomrule
\end{tabular}

\begin{apcnote}[Scoring yourself honestly]
12/12: move on to the unit worksheets and the free-response set.
9--11: solid --- redo the missed ladders tomorrow from a blank page.
8 or fewer: the recurring root causes here are exactly three ---
(1) forgetting the coefficient as a multiplier of $x$ in an ICE table,
(2) claiming $K$ changed when only the position moved, and (3) writing
$K_{sp} = s^n$ without the factor that the stoichiometry demands. Fix
those three and most of these ladders fall together.
\end{apcnote}

\end{document}
